\scanpage{46}
\mspage{44}
% archive sheet number: 45
9. \uline{Réductions au cas projectif et lisse.}

9.0. Nous nous posons la question de savoir, pour un \inserted{$V$ et} un entier $\unsure{n|\ell}$ donné, si le th. de monodromie 7.1. est vrai pour tout $X_K$ de t.f. sur $K$, \deleted{de dim} \inserted{tel que}
\[
\dim X_K\leq n.
\]
Nous allons d'ailleurs envisager la conjecture correspondante également pour les $\pi$-modules $H^i_!(X_{\overline K},\deleted{\unsure{\mathbb Z|Z}}\unsure{\mathbb Q_\ell|Q_\ell})$ (pour $X_K$ \uline{séparé}).
\inserted{Nous supposons que la conjecture est prouvée déjà pour les dim. $\leq n-1$.}

9.1. \uline{Cas $X$ lisse.} \inserted{\uline{Relation entre $H^i$ et $H^i_!$}} Dans ce cas, \deleted{\unsure{et}} \inserted{et si $X$ est séparé,} on a un isom canonique
\[
H^i(X_{\overline K},\deleted{\unsure{\mathbb Z}}\mathbb Q_\ell)\simeq H^{\deleted{\unsure{i}}\inserted{2n-i}}_!(X_{\overline K},\mathbb Q_\ell)'\deleted{\unsure{\otimes_{\mathbb Q_\ell}\mathbb Q_\ell}}(-n),
\]
où \deleted{$\unsure{\otimes}$} $(-n)$ désigne le twist de Tate $\otimes_{\mathbb Q_\ell}\mathbb Q_\ell(-n)$, \deleted{\unsure{et}} et $'$ désigne le dual. Cet isom. est compatible avec l'action de $\pi$ et à fortiori de $I$. Comme $I$ opère trivialement sur $\mathbb Q_\ell(-n)$, il s'ensuit que l'action de $I$ sur $H^i(X_{\overline K};\mathbb Q_\ell)$ est quasi unip. (resp. unipot.) ssi son action sur $H^i_{\deleted{\unsure{*}}!}(X_{\overline K},\mathbb Q_\ell)$ l'est.

9.2. \deleted{\unsure{Le cas $X_K$ séparé}} \inserted{\uline{Cas} $H^i_!(X_{\overline K},\mathbb Q_\ell)$, $X_K$ séparé de t.f. sur $K$.} \deleted{Pour} Soit $U$ un ouvert \deleted{affine} \inserted{dense} de $X_K$, $Y$ son complémentaire. Alors on a la suite exacte
\[
\cdots\longrightarrow H^i_!(\deleted{\unsure{U}}\unsure{Y_{\overline K}|U_{\overline K}},\deleted{\unsure{\mathbb Z_\ell|\mathbb Q_\ell}}\mathbb Q_\ell)\longrightarrow H^i_!(X_{\overline K},\mathbb Q_\ell)-H^i_!(U_{\overline K},\mathbb Q_\ell)-H^{i+1}_!(U_{\overline K},
\]
% The right edge of the sheet cuts the last displayed term after the comma.
qui montre, moyennant l'hyp. de récurrence, que
